\section{Simulating Human Subjects: Behavioural Fidelity and the Hyper-Accuracy Distortion}
\label{sec:persona_simulation}

\subsection{The Turing Experiment Framework}
Aher, Arriaga and Kalai~\cite{aher2023using} introduce the \emph{Turing Experiment} (TE). Unlike the Turing Test, which asks whether a machine can pass as one individual, a TE asks whether a language model can simulate a \emph{representative sample} of participants in a human-subject study. Participants are represented by names with gender titles (e.g.\ ``Ms.\ Olson''), and the model completes a prompt that describes the experimental situation. To reduce the risk that the model simply recalls published results, the authors wrote new material for three of the four studies, for example their own garden-path sentences and their own obedience scenario~\cite{aher2023using}. % SOURCE: Aher et al. p.2 (names/titles), p.3 and p.7 (own variations; 24 own garden-path sentences, own obedience scenario)

\subsection{Replication Across Behavioural Paradigms}
The authors ran four TEs with GPT models accessed through the OpenAI API; in the first three, larger models gave more faithful simulations and the largest ones replicated the known human findings~\cite{aher2023using}: % SOURCE: Aher et al. p.1 abstract; p.3 (larger models more faithful); p.7 (model list)
\begin{enumerate}
    \item \textbf{Ultimatum Game (behavioural economics):} with \texttt{text-davinci-002}, simulated responders almost always accepted offers of 50--100\% of the endowment and rarely accepted offers of 0--10\%, in line with human behaviour; smaller models were not sensitive to the offer size. % SOURCE: Aher et al. p.8 (LM-5); LM-5 = text-davinci-002 per p.7
    \item \textbf{Garden-path sentences (psycholinguistics):} for locally ambiguous sentences such as \emph{``While the student read the notes that were long and boring blew off the desk''}, the largest models rated garden-path sentences as ungrammatical more often than control sentences, as expected from human studies. % SOURCE: example p.2 Fig. 1; hypothesis p.19; largest models replicate the first three TEs p.3
    \item \textbf{Destructive obedience (social psychology):} in Milgram's original study, 26 of 40 participants followed the instructions to the end of the shock series~\cite{milgram1963behavioral}; in the simulated version with the authors' new scenario, 75 of 100 simulated participants did so. % SOURCE: Aher et al. p.11 (26/40 human, 75/100 simulated)
\end{enumerate}

\subsection{The Hyper-Accuracy Distortion}
The fourth TE concerns the \emph{wisdom of crowds}: individual human estimates of a quantity are noisy, but their aggregate is often close to the truth~\cite{surowiecki2004wisdom}. Here the simulation fails in a characteristic way. For the question ``How many bones does an adult human have?'', human participants gave a median answer of 190 with an interquartile range (IQR) of 108, whereas \texttt{gpt-3.5-turbo} and GPT-4 gave the exact answer of 206 with an IQR of 0; the older \texttt{davinci} model gave a median of 136 (IQR 346). % SOURCE: Aher et al. Table 3 p.25 (human median 190, IQR 108; truth 206) and Table 4 p.26 (davinci 136/346; gpt-35-turbo 206/0; gpt-4 206/0)
The same pattern appears for other questions, e.g.\ the melting temperature of aluminium. Aher et al.\ call this the \emph{hyper-accuracy distortion} and observe that it is stronger in the more recent, more aligned models; they suggest that it \emph{may} be caused by alignment procedures that improve truthfulness~\cite{aher2023using, ouyang2022training}. % SOURCE: Aher et al. p.3 (aluminium example, "may be due to alignment"), Fig. 15 caption p.26
For simulation studies this means that aligned models can reproduce qualitative behavioural patterns but cannot be used, without adjustment, to simulate the natural spread of human knowledge and error.
